% CV of Xiaoyang Cao. Build with:  xelatex CV_Xiaoyang_Cao.tex
% then copy the PDF to ../CV_Xiaoyang_Cao.pdf and ../CV_Xiaoyang_Cao_2026.pdf.
\documentclass[10pt,letterpaper]{article}

\usepackage[margin=0.78in,top=0.62in,bottom=0.7in]{geometry}
\usepackage{fontspec}
\setmainfont{Latin Modern Roman}
\usepackage{enumitem}
\usepackage{titlesec}
\usepackage[hidelinks]{hyperref}

\pagestyle{empty}
\setlength{\parindent}{0pt}
\setlength{\parskip}{0pt}
\linespread{1.0}

% Section headings: bold caps with a full-width rule.
\titleformat{\section}{\large\bfseries}{}{0pt}{\MakeUppercase}[\vspace{2pt}\titlerule]
\titlespacing*{\section}{0pt}{12pt}{12pt}

% Links in the body are underlined; [arXiv]-style tags are bold.
\newcommand{\ulink}[2]{\href{#1}{\underline{#2}}}
\newcommand{\linktag}[2]{\textbf{[\ulink{#1}{#2}]}}
\newcommand{\sep}{\enspace$\cdot$\enspace}

% \entry{bold title}{right}{italic subtitle}{right}
\newcommand{\entry}[4]{%
  \textbf{#1}\hfill #2\\
  \textit{#3}\hfill #4\par}

% \project{title}{dates}{advisor line}
\newcommand{\project}[3]{%
  \vspace{6pt}%
  {#1}\hfill #2\\
  \textit{#3}\par\vspace{2pt}}

\setlist[itemize]{leftmargin=12pt,labelsep=5.5pt,itemsep=1.4pt,parsep=0pt,topsep=1pt,label=\textbullet}
\setlist[enumerate]{leftmargin=*,labelindent=0pt,labelsep=5pt,itemsep=7.4pt,parsep=0pt,topsep=0pt,label={[\arabic*]},widest=5}

\begin{document}

\begin{center}
  {\LARGE\bfseries Xiaoyang Cao}\\[5pt]
  Cambridge, MA\sep (+1) 617-216-8184\sep \href{mailto:xycao@mit.edu}{xycao@mit.edu}\sep
  \href{https://xiaoyangcao1113.github.io/}{Personal Website}\sep
  \href{https://scholar.google.com/citations?user=XVWuYHAAAAAJ}{Google Scholar}\sep
  \href{https://github.com/XiaoyangCao1113}{GitHub}
\end{center}
\vspace{-6pt}

\section{Education}

\entry{Massachusetts Institute of Technology}{Cambridge, MA}
      {M.S. in Computational Science \& Engineering}{Sep 2025 -- May 2027 (expected)}
\textit{Concurrent M.S. in Technology \& Policy}\\
\textit{Laboratory for Information and Decision Systems (LIDS)}\\
\textit{Research advisor: Prof. \ulink{https://miba.dev/}{Michiel Bakker}}\par
\vspace{6pt}
\entry{Tsinghua University}{Beijing, China}
      {B.S. in Mathematics and Physics}{Sep 2021 -- Jun 2025}
\vspace{6pt}
\entry{University of California, Berkeley}{Berkeley, CA}
      {Exchange Student}{Jan 2024 -- Aug 2024}

\section{Research Interests}

How intelligence emerges in collectives of interacting agents (multi-agent reinforcement learning, populations of
LLM agents) and in self-referential learning (self-play, self-distillation, preference optimization).

\section{Publications}

\begin{enumerate}
  \item \textbf{Xiaoyang Cao}, Jingqi Li, Zhe Fu, Alexandre M. Bayen. ``Who Bears the Burden? Learning Responsibility for
    Shared Constraints in Multi-Agent Reinforcement Learning.'' \textit{Under review}, 2026.
    \href{https://arxiv.org/abs/2610.07491}{arXiv:2610.07491}.
  \item \textbf{Xiaoyang Cao}, Siddarth Srinivasan, Michiel A. Bakker. ``Do Modules Stay in Their Lane? Role Drift in
    Compound LLM Systems.'' \textit{Under review}, 2026.
    \href{https://arxiv.org/abs/2607.21627}{arXiv:2607.21627}.
  \item \textbf{Xiaoyang Cao}, Zelai Xu, Mo Guang, Kaiwen Long, Michiel Bakker, Yu Wang, Chao Yu. ``RE-PO: Robust
    Enhanced Policy Optimization as a General Framework for LLM Alignment.'' \textit{International Conference on
    Learning Representations (ICLR)}, 2026.
  \item \textbf{Xiaoyang Cao}, Zhe Fu, Alexandre M. Bayen. ``Pareto Control Barrier Function for Inner Safe Set Maximization
    Under Input Constraints.'' \textit{American Control Conference (ACC)}, 2025.
  \item \textbf{Xiaoyang Cao}, Dingyi Zhuang, Shenhao Wang, Jinhua Zhao. ``Virtual Nodes Improve Long-term Traffic
    Prediction.'' \textit{Transportation Research Board Annual Meeting (TRB)}, 2025.
\end{enumerate}

\section{Research Experience}

\vspace{-6pt}
\project{\textbf{Social Swarm: Scaling Populations of Research Agents with Sparse Networks}}{Sep 2026 -- Present}
        {Advised by Prof. \ulink{https://miba.dev/}{Michiel Bakker}, MIT}
\begin{itemize}
  \item Studying how collectives of LLM research agents convert agent count into capability on open-ended problems,
    motivated by recent results that scale to thousands of agents.
  \item Identified \textit{diversity collapse}: when every agent sees every attempt, the population homogenizes and long-horizon
    performance stalls. Proposed a sparse follow-network organization in which each agent sees only a bounded set
    of peers.
  \item Preliminary experiments on the Heilbronn triangle problem show that the sparse network delays collapse relative
    to a fully visible shared board and to independent best-of-$N$.
\end{itemize}

\project{\textbf{Learning Responsibility for Shared Constraints in MARL} \textbf{[1]}
         \linktag{https://arxiv.org/abs/2610.07491}{arXiv} \linktag{https://lira-marl.github.io/}{Project}}{May 2025 -- Sep 2026}
        {Advised by Prof. \ulink{https://bayen.berkeley.edu/alex-bayen}{Alexandre Bayen}, UC Berkeley}
\begin{itemize}
  \item Proposed LiRA (Lagrangian Responsibility Allocation): a common multiplier enforces a shared cost budget, while
    learned per-agent responsibility shares redistribute its penalty by optimizing social welfare, without modifying
    rewards or constraints.
  \item Proved that for convex games varying the shares induces a smooth family of normalized generalized Nash equilibria
    in which active constraints stay at budget while welfare varies, and derived a welfare gradient that accounts for
    both learning updates and the induced shift in data distribution.
  \item Improved average social welfare by up to 29\% over uniform and agent-specific multiplier baselines across CityLearn,
    MABIM, Harvest, and MetaDrive (3 to 400 agents).
\end{itemize}

\project{\textbf{Role Drift in Compound LLM Systems} \textbf{[2]}
         \linktag{https://arxiv.org/abs/2607.21627}{arXiv}
         \linktag{https://venturebeat.com/ai/one-ai-module-faked-86-of-a-pipelines-accuracy-gains-by-feeding-another-the-answers}{VentureBeat}}{Apr 2026 -- Sep 2026}
        {Advised by Prof. \ulink{https://miba.dev/}{Michiel Bakker}, MIT}
\begin{itemize}
  \item Identified \textit{role drift}: end-to-end reinforcement learning can improve terminal performance while causing specialized
    modules to abandon their intended functions.
  \item Developed role-level evaluations showing that 86\% of an apparent 31-percentage-point accuracy gain disappeared
    once the decomposer's assigned function was measured, demonstrating that end-to-end metrics can conceal
    internal failure.
  \item Proposed Role Anchor, a training objective that preserves the influence of role specifications while allowing
    modules to improve under end-to-end optimization.
\end{itemize}

\project{\textbf{RE-PO: Robust LLM Preference Alignment} (ICLR 2026) \textbf{[3]}
         \linktag{https://arxiv.org/abs/2509.24159}{arXiv} \linktag{https://repo-alignment.github.io/}{Project}}{Mar 2025 -- Sep 2025}
        {Advised by Prof. \ulink{https://web.ee.tsinghua.edu.cn/wangyu/en/}{Yu Wang}, Tsinghua University}
\begin{itemize}
  \item Proposed an EM-based framework that models preference-label correctness as a latent variable, jointly estimates
    label and annotator reliability, and systematically robustifies DPO, IPO, SimPO, and CPO.
  \item Proved convergence to the true noise level under calibration and improved AlpacaEval 2 win rates by up to 7.0
    percentage points on Mistral-7B and Llama-3-8B, with additional gains on Arena-Hard and MultiPref.
\end{itemize}

\project{\textbf{Safe Control via Pareto Control Barrier Functions} (ACC 2025) \textbf{[4]}
         \linktag{https://arxiv.org/abs/2410.04260}{arXiv}}{May 2024 -- Oct 2024}
        {Advised by Prof. \ulink{https://bayen.berkeley.edu/alex-bayen}{Alexandre Bayen}, UC Berkeley}
\begin{itemize}
  \item Developed the Pareto Control Barrier Function algorithm by integrating Pareto multi-task learning into neural
    control barrier functions, maximizing the inner safe set while guaranteeing safety under input constraints.
  \item Matched the Hamilton--Jacobi viability kernel on a 2D inverted pendulum and achieved a 700\% larger inner safe
    set than the benchmark on a 12-D quadrotor obstacle-avoidance task.
\end{itemize}

\project{\textbf{Virtual Nodes for Long-term Traffic Prediction} (TRB 2025) \textbf{[5]}
         \linktag{https://arxiv.org/abs/2501.10048}{arXiv}}{May 2024 -- Sep 2024}
        {Advised by Prof. \ulink{https://zhaojinhua.com/}{Jinhua Zhao}, MIT}
\begin{itemize}
  \item Introduced virtual nodes into spatio-temporal graph neural networks to mitigate over-squashing in long-horizon
    traffic forecasting, using a semi-adaptive adjacency matrix learned from node embeddings.
  \item Reduced RMSE by 6.27\% and MAPE by 5.04\% on a five-year, 716-sensor San Diego dataset; the learned adjacency
    emphasized traffic-intensive intersections, improving interpretability.
\end{itemize}

\section{Awards \& Honors}

\textbf{Outstanding Winner \& INFORMS Award}, Mathematical Contest in Modeling -- top 1\%\hfill May 2023\\
\textbf{Outstanding Innovation in Science \& Technology Scholarship}, Tsinghua University\hfill 2023, 2024\\
\textbf{Comprehensive Excellence Scholarship}, Tsinghua University\hfill Oct 2023\\
\textbf{Outstanding Academic Scholarship}, Tsinghua University\hfill Oct 2022

\section{Service \& Teaching}

\textbf{Conference Reviewer:} ICLR 2027, CoRL 2026\\
\textbf{Teaching Assistant:} Calculus A, Tsinghua University\hfill Spring 2025

\section{Skills \& Languages}

\textbf{ML \& RL:} Python, PyTorch, Hugging Face Transformers / TRL, vLLM, DeepSpeed, FSDP, PettingZoo, RLlib\\
\textbf{Compute:} Multi-GPU, multi-node training on H100 / H200 / B200 clusters\\
\textbf{Languages:} English (fluent), Chinese (native)

\end{document}
